\documentclass[10pt]{article}

\usepackage[utf8]{inputenc}

\usepackage[T1]{fontenc}

\usepackage{amsmath}

\usepackage{amsfonts}

\usepackage{amssymb}

\usepackage[version=4]{mhchem}

\usepackage{stmaryrd}

\usepackage{bbold}



\begin{document}

речь идет о качественной картине поведения всех траекторий во всем фазовом пространстве (глобальная теория) или по крайней мере в нек-рой его части (локальная теория); в теории Д. с. особенно большое внимание уделяется поведению фазовых траекторий при неограниченном возрастании времени. Из отдельных же траекторий в теории Д.с. интересуются обычно теми, свойства к-рых могут в значительной степени влиять на качественную картину, хотя бы локальную.



Для системы двух уравнений вида (1) $(m=2)$ кинематич. интерпретация дает весьма наглядный и эффективный способ исследования, поскольку векторное поле $f(w)$ и фазовые траектории могут быть практически изображены на фазовой плоскости. Уже в случае трех уравнений $(m=3)$ пришлось бы производить соответствующие построения в 3-мерном пространстве, что довольно затруднительно, а при $m>3$ подобные практич. возможности вообще отпадают. Поэтому при $m \geqslant 3$ (а в ряде случаев и при $m=2$ ) роль кинематич. интерпретации состоит в использовании при исследовании дифференциальных уравнений геометрич. понятий, методов и языка, в той или иной степени обобщающих привычные из повседневной жизни геометрич. представления.



Уже сравнительно слабые предположения о векторном поле $f(w)$ (напр., его дифференцируемость) гарантирует то, что для каждой точки $w_{0} \in \boldsymbol{W}^{\boldsymbol{m}}$ существует ровно одно решение $w(t)$ уравнения (3), имеющее $w_{0}$ своим начальным значением: $w(0)=w_{0}$. Физически это соответствует тому, что при заданном законе движения (3) состояние системы в любой момент времени полностью определяется ее начальным состоянием. Вообще говоря, это решение может быть определено не для всех $t$, а лишь на нек-ром отрезке времени. В глобальной теории Д. с. делается дополнительное предположение, что при любом начальном значении соответствующее решение определено при всех $t$, тогда как в локальных вопросах обычно нецелесообразно делать какие-либо предположения о дальнейшем поведении тех траекторий, к-рые покидают рассматриваемую область фазового пространства.



При выполнении указанного предположения, если каждому $w_{0} \in W^{m}$ сопоставить то состояние, куда перейдет за время $t$ фазовая точка, движущаяся согласно (3) и вышедшая при $t=0$ из $w_{0}$, то получится нек-рое отображение $S_{t}$ фазового пространства $W^{m}$ в себя:



\[

S_{t} w_{0}=w(t)

\]



где $w(t)-$ решение (3) и $w(0)=w_{0}$. Отображения $S_{t}$ образуют непрерывную однопараметрич. группу диффеоморфизмов фазового многообразия $W^{m}$ [групповое свойство $S_{t} S_{s}=S_{t+s}$ является следствием автономности системы (3)]. В порядке иллюстрации в литературе нередко проводят аналогию с известным из повседневной жизни и ранее всего изученным в науке примером, где возникает подобное семейство преобразований пространства,- стационарным течением жидкости или газа: за время $t$ частица жидкости перетекает из точки $w_{0}$ в $S_{t} w_{0}$. (Впрочем, эта аналогия довольно поверхностна, ибо воображаемая «фазовая жидкость», «текущая» в фазовом пространстве, отличается от реальных сплошных сред отсутствием взаимодействия между соседними частицами.) В связи с этим в качестве синонима термина «Д.с.» употребляется термин «П о то К».



В физич. литературе принято говорить об ан с ам бле д и на м и ч е ск и х с и стем. Это означает, что каждому из возможных состояний данной физич. системы (то есть Каждой точке фазового пространства) мысленно сопоставляется нек-рая физич. система, описываемая уравнением (3) и находящаяся в этом состоянии; полученную совокупность однотипных систем (никак не взаимодействующих друг с другом), различающихся только состояниями, к-рыми они обладают в данный момент, и называют «ансамблем». На



176 ДИНАМИЧЕСКАЯ этом языке преобразованиям $S_{t}$ фазового пространства соответствует эволюция «ансамбля», заключающаяся в изменении состояния входящих в него систем.



При разработке глобальных аспектов теории Д. с. понятие Д.с. подверглось дальнейшему обобщению. В наиболее широком смысле под д и на м и ч е ской с и с те м о й понимают произвольное действие группы (или даже полугруппы) $G$ на нек-ром множестве $W$, именуемом фазовым пространством. Это значит, что для любого $g \in G$ определено отображение $S_{g}: W \rightarrow W$, причем $S_{g_{1}} S_{g_{2}}=S_{g_{1} g_{2}}$, и если $e-$ единица $G$, то $S_{e}$ - тождественное преобразование (то есть $S_{e}(w)=w$ для всех $\left.w\right)$. Совокупность точек $S_{g}\left(w_{0}\right)$, где $w_{0}$ фиксировано, a $g$ пробегает $G$, называется т ра е к тор и ей (или о р б и т о й), проходящей через точку $w_{0}$, или, короче, траекторией этой точки. Группа $G$ обычно считается топологич. группой, фазовое пространство - топологич. пространством, или пространством с мерой, а отображение



\[

G \times W \rightarrow W,(g, w) \rightarrow S_{g}(w)

\]



предполагается, соответственно, непрерывным или измеримым, причем в последнем случае обычно предполагается, что отображения $S_{g}$ сохраняют меру (то есть прообраз измеримого подмножества фазового пространства измерим и имеет ту же меру). Соответствующие направления в теории Д.с. называются топол о гичес ской динам икой (cм. [6], [11]) и эргодической тео рией (см. [4], [5], [7], [12]). Если $G$ - группа Ли, $W$ - гладкое многообразие, а отображение (5) гладкое, то говорят о глад кой ди намической системе.



Впрочем, во всех трех направлениях основными являются случаи, когда $G$ - либо группа $\mathbb{R}$ действительных чисел, и тогда Д. с. называют пото к о м (хотя иногда этот термин употребляют как синоним термина «Д. с.» в наиболее общем его значении), либо $G$ - группа $\mathbb{Z}$ целых чисел, и тогда для таких Д.с. был предложен термин каскад (говорят также о динамической системе с дискретным в р е м е н е м, но этот термин может означать также и только то, что в $G$ берется дискретная топология). Гладкий поток, для к-рого отображение (5) - класса $C^{2}$, определяется гладким векторным полем



\[

f(w)=\left.\frac{d}{d t} S_{t}(w)\right|_{t=0}

\]



\begin{itemize}

  \item при изменении $t$ точка $w=S_{t}\left(w_{0}\right)$ движется согласно (3). В случае каскада отображения $S_{n}$ получаются итерированием преобразования $S_{1}$ и обратного к нему (вариант: итерированием одного только отображения $S_{1}$ ); для гладкого каскада все $S_{n}$ суть диффеоморфизмы.

\end{itemize}



Среди названных трех направлений в теории Д. с. топологич. динамика имеет ярко выраженный теоретикомножественный характер, а ее роль вначале была как бы вспомогательной. Дело в том, что обсуждение ряда понятий (неблуждающая точка, предельное множество траекторий, минимальное множество, почти периодичность, дистальность, устойчивость по Лагранжу, устойчивость по Пуассону и др.) и связей между ними, важных для трактовки более конкретных объектов - гладких Д.с.,- целесообразно провести в более абстрактной обстановке, отвлекаясь от задания Д. с. при помощи диффеоморфизма или уравнения (3). Это и делается в топологич. динамике. Позднее она получила значительное развитие, главным образом, в направлении изучения нек-рых минимальных множеств и их расширений (см. [13]-[15]).



Зарождение эргодич. теории связано с классич. (доквантовой) статистич. физикой. При обосновании последней возник вопрос, можно ли, не решая гамильтонову систему дифференциальных уравнений, к-рая описывает движение частиц, образующих рассматриваемое макроскопич. тело, и даже не зная начальные значения для ее решения (что означало бы задание мгновенных положений и скоростей всех этих частиц), указать свойства статистич. характера, проявляюшиеся в поведении при $t \rightarrow \infty$ всех или почти всех





\end{document}